\chapter{总结与展望}

\section{全文工作总结}

本文以同济大学 Unix V6++ 教学操作系统为对象，在团队整体 Rust 重构和 RISC-V 迁移工作的基础上，重点研究内存管理、文件系统及其相关缓冲区与块设备路径。课题的基本思路不是脱离原系统重新设计教学内核，而是在保留 Unix V6++ 课程结构和主要算法的前提下，重新整理资源生命周期、错误传播路径和平台相关代码边界。

在原系统分析阶段，本文梳理了 Unix V6++ 的整体结构及其 IA32 平台依赖。原系统以 C++ 实现，覆盖进程、内存、文件和设备等基本模块，具有较好的教学可读性。但其页表、inode、缓冲区和打开文件对象等共享资源主要依赖裸指针、全局状态和手工约定维护；IA32 分页、中断控制器和 ATA 设备接口也使底层代码与具体平台耦合较深。这些问题不会否定原系统的教学价值，却会增加后续扩展、调试和跨架构迁移的成本。

从分支差异看，\texttt{main} 分支保留了原始 C++/IA32 结构，内存管理、文件系统、缓冲区缓存、设备驱动、进程管理和中断处理均以 C++ 模块组织。\texttt{refactor/x86} 分支将内核主体切换为 Rust 实现，并在 IA32 平台上保留 ATA、DMA、I/O 端口、8259A 和 x86 页表等平台路径。\texttt{refactor/riscv} 分支则在 Rust 版本基础上进一步替换体系结构相关代码，删除 ATA、DMA 和 I/O 端口路径，引入 VirtIO 块设备、PLIC 外部中断、RISC-V trap 处理和 Sv39 页表适配。本文的分析重点放在这三类变化对内存管理、文件系统及其 I/O 路径的影响上。

在内存管理模块中，Rust 版本重新组织了物理页、页级分配、内核堆和交换区管理。物理页以统一结构表示，页框号、物理地址和地址区间通过类型化抽象区分；Buddy 分配器负责页级连续空间分配，Slab 分配器负责内核小对象分配，两者共同改善了原系统分配粒度较粗和对象边界不够明确的问题。交换区管理以显式空闲区间表组织分配、释放和合并逻辑，减少了空间管理代码在过程式路径中的分散程度。

在文件系统模块中，Rust 版本保留了 Unix V6 风格的超级块、inode 缓存、目录项查找、打开文件表和多级块映射机制，重构重点放在对象引用关系和错误处理路径上。内存 inode、文件对象和文件描述符表通过包装类型表达共享和借用关系，资源释放依赖析构路径完成；内部错误使用 \texttt{Result<T, E>} 与 \texttt{PosixError} 组织，并在系统调用边界转换为符合 Unix 语义的错误返回。缓冲区缓存和设备层则通过侵入式链表、状态字段和 trait 抽象重新组织，使文件系统与块设备之间的接口更清晰。

在 RISC-V 迁移部分，本文围绕内存管理和文件系统所依赖的底层条件，分析并适配了 Sv39 分页、UART 串口和 VirtIO 块设备等路径。RISC-V 版本重新建立虚拟地址空间布局，将内核映射、用户空间和 MMIO 设备区域区分开来；块设备路径从 IA32 版本的 ATA/DMA 机制切换到 VirtIO MMIO 接口。文件系统核心结构在迁移后基本保持不变，说明重构后的上层文件系统逻辑能够与平台相关设备路径相分离。

测试部分对原始 C++ IA32、Rust IA32 和 Rust RISC-V 三个版本进行了比较。功能测试表明，两个 Rust 版本均能完成文件系统装载、控制台初始化、shell 运行、文件读写和用户程序执行等基本任务。性能测试覆盖系统调用延迟、用户态动态分配和文件 I/O 吞吐量。结果显示，Rust 重构没有引入不可接受的整体开销，部分系统调用路径在重构后还表现出更低延迟；RISC-V 版本在小规模系统调用测试中表现接近 IA32 版本，但在较大文件 I/O 场景下仍受到 VirtIO 驱动实现和设备路径的影响。

\section{主要成果}

本文的主要成果体现在四个方面。

第一，完成了 Unix V6++ 内存管理模块的 Rust 重构。重构后的实现以类型化地址、物理页包装、Buddy 分配器、Slab 分配器和交换区空闲区间表组织内存资源，减少了裸地址和裸指针在上层逻辑中的直接传播。RAII 风格的资源回收使页对象和内核对象的生命周期更容易从代码结构中识别。

第二，完成了 Unix V6++ 文件系统及缓冲区缓存路径的 Rust 重构。重构后的文件系统保留了超级块、inode、目录项、打开文件表和多级块映射等教学核心内容，同时通过引用封装、自动析构和显式错误类型整理了 inode、文件对象和缓冲区之间的资源关系。块设备和字符设备以 trait 表达统一接口，便于后续替换底层驱动。

第三，完成了与内存管理和文件系统相关的 RISC-V 迁移适配。系统已建立 Sv39 页表和新的地址空间布局，块设备路径由 ATA/DMA 转为 VirtIO MMIO，文件系统能够在 RISC-V 平台完成装载并支撑 shell 中的文件访问。该结果说明 Unix V6++ 的教学结构可以在新的体系结构上继续工作，也说明将平台相关代码集中到底层模块有助于跨架构适配。

第四，形成了三版本对比测试结果。本文从功能正确性和基础性能两个维度比较了 C++ IA32、Rust IA32 和 Rust RISC-V 三个版本。测试结果为后续课程实验提供了基准数据，也揭示了不同实现路径在系统调用、内存分配和文件 I/O 场景下的差异。

\section{存在的不足}

本文仍存在若干不足。研究边界方面，Unix V6++ 的 Rust 重构和 RISC-V 迁移是团队协作完成的系统工程，本文重点讨论内存管理、文件系统以及二者依赖的缓冲区和块设备路径。进程管理、系统调用分发、中断处理和终端等模块虽已在相关分支中参与系统集成，但其设计细节属于其他课题的主要工作范围。因此，本文对这些模块只在接口依赖和测试现象层面展开分析，未对其内部实现作同等深度的评价。

内存管理方面，当前交换区采用较直接的空闲区间管理策略，在频繁换入换出场景下可能产生碎片。Buddy 分配器阶数和 Slab 槽位范围主要由静态常量确定，缺少运行时配置能力。RISC-V 版本目前以 4KB 页为基本映射单位，尚未利用大页减少内核映射中的页表层数。

文件系统方面，当前实现保留 Unix V6 风格的固定大小目录项和有限 inode 地址表，文件大小和目录组织能力受到原始结构限制。系统暂不支持多挂载点，也没有实现更复杂的文件系统一致性恢复机制。该取舍符合教学系统的简化目标，但限制了其作为通用实验平台的扩展空间。

设备与驱动方面，本文关注的是文件系统下层块访问路径在 IA32 ATA/DMA 与 RISC-V VirtIO MMIO 之间的适配关系。当前 RISC-V 版本的 VirtIO 块设备路径能够满足文件系统装载和基本读写需求，但尚未围绕请求合并、异步 I/O 和中断驱动传输进行充分优化。第 7 章中 RISC-V 版本在较大文件 I/O 测试中的差距，主要应理解为设备路径和模拟平台共同作用的结果，而不是文件系统上层算法本身发生了根本变化。

测试方面，当前测试覆盖了基本功能路径和若干微基准，但针对内存分配边界、inode 引用释放、目录异常路径、缓冲区缓存替换和块设备错误返回的系统化测试仍不充分。性能测试受 QEMU 环境、系统时钟粒度和宿主机负载影响，结果适合用于版本间相对比较，不宜解释为硬件平台上的绝对性能结论。unsafe 边界方面，本文已尽量在内存映射、页表、MMIO 和底层设备访问等路径中明确其作用范围，但尚未形成系统化的 unsafe 审查文档。

\section{后续工作展望}

后续工作可以从测试验证、性能优化和教学支撑三个方向推进。

在测试验证方面，可围绕本文涉及的内存管理和文件系统路径补充更细粒度的单元测试与集成测试。内存管理可增加页分配回收、Slab 对象复用、交换区碎片合并和页表权限检查测试；文件系统可增加 inode 引用计数、目录项边界、文件截断、错误返回和缓冲区缓存一致性测试。这类测试比单纯增加用户态功能程序更能暴露资源生命周期错误。

在内存和文件系统方面，可为 Slab 分配器增加多尺寸对象缓存和按对象类型划分的缓存链表；交换区可改用位图或更细粒度的空闲页管理策略；RISC-V 版本可探索大页映射以降低页表开销。文件系统后续可扩展 inode 地址表、支持更大文件和多挂载点，并补充一致性检查或只读挂载等机制。

在设备和性能方面，VirtIO 块设备驱动可增加请求合并和中断驱动 I/O 路径，减少轮询等待对吞吐量的影响。性能测试可引入 RISC-V cycle 计数器或 QEMU 支持的更高精度计时手段，以降低 60Hz tick 粒度对结果解释的影响。对于跨架构比较，还需要固定宿主机负载、QEMU 版本、虚拟机内存和编译优化选项，以减少实验环境差异对结果的干扰。

在教学支撑方面，可围绕重构后的内存管理和文件系统补充实验文档、代码阅读路线和 unsafe 边界说明。对于教学系统而言，Rust 重构的价值在于减少部分内存错误，并把资源所有权、错误处理和平台适配边界明确呈现给学生。后续若能将这些内容转化为实验任务，Unix V6++ 可以继续作为理解传统 Unix 内核和现代系统软件实现方法之间关系的课程平台。
